% SPDX-License-Identifier: CC-BY-4.0
% Copyright 2025 Florian Aram Feuerriegel — kassensturz.org
% Anhang — Daten, Ausgaben, Messwerte, Ressortgraphen, Elastizitäten, Quellen
% (alle Tabellen und Graphen erzeugt von messung.js). Im Anhang gleitet nichts: Jede Tabelle
% und jeder Graph steht unter seiner Überschrift, \captionof statt table/figure.

\newenvironment{anhangblock}{\par\medskip\noindent\begin{minipage}{\textwidth}\centering}%
  {\end{minipage}\par\medskip}

\section{Haushaltsdaten}
\label{sec:anhang-daten}

\begin{anhangblock}
\captionof{table}{Die zwölf Haushaltstypen (\datei{DEZILE} und Profile in
\datei{js/data.js}). $B_i$ Bruttoeinkommen, $\kappa_i$ Kapitalanteil, $c_i$
durchschnittliche Konsumquote, $w_i$ Bedarfsgewicht (neue OECD-Skala), $V_i$
Vermögen, $N_i$ Haushalte, $\rho_i$ Rentenanteil am Einkommen (Näherung),
$m_i$ Grenzkonsumneigung (Näherung), $k_i$ Kindergeldkinder je Haushalt, $q_i$
Bürgergeld-Regelsatzjahre je Haushalt, $\ell_i$ \COz-Last je Haushalt im
Status quo (55\,\euro/t, vollständig überwälzt).}
\label{tab:dezile}
{\footnotesize\setlength{\tabcolsep}{3.5pt}% Erzeugt von docs/modell/messung.js — nicht von Hand bearbeiten
\begin{tabular}{@{}lrrrrrrrrrrr@{}}
\toprule
Typ & $B_i$ & $\kappa_i$ & $c_i$ & $w_i$ & $V_i$ & $N_i$ & $\rho_i$ & $m_i$ & $k_i$ & $q_i$ & $\ell_i$ \\
 & Tsd.\,\euro{} & \% & \% & & Tsd.\,\euro{} & Mio. & \% & & & & \euro{} \\
\midrule
D1 & 14 & 1 & 100 & 1,3 & 1 & 4,10 & 42,0 & 0,65 & 0,37 & 0,60 & 166 \\
D2 & 21 & 1 & 99 & 1,4 & 5 & 4,10 & 46,2 & 0,62 & 0,51 & 0,25 & 236 \\
D3 & 27 & 2 & 96 & 1,5 & 15 & 4,10 & 42,0 & 0,58 & 0,56 & 0,08 & 287 \\
D4 & 33 & 2 & 92 & 1,6 & 35 & 4,10 & 35,3 & 0,55 & 0,54 & 0,02 & 332 \\
D5 & 40 & 3 & 88 & 1,7 & 70 & 4,10 & 28,6 & 0,52 & 0,49 & 0,00 & 379 \\
D6 & 48 & 3 & 85 & 1,8 & 120 & 4,10 & 21,8 & 0,48 & 0,44 & 0,00 & 426 \\
D7 & 58 & 4 & 82 & 1,9 & 200 & 4,10 & 15,1 & 0,45 & 0,40 & 0,00 & 480 \\
D8 & 72 & 5 & 78 & 2,0 & 340 & 4,10 & 10,1 & 0,42 & 0,35 & 0,00 & 532 \\
D9 & 95 & 7 & 72 & 2,0 & 620 & 4,10 & 6,7 & 0,38 & 0,30 & 0,00 & 618 \\
D10a & 125 & 8 & 60 & 2,0 & 850 & 2,05 & 4,2 & 0,32 & 0,23 & 0,00 & 665 \\
D10b & 220 & 18 & 52 & 2,0 & 1.500 & 1,64 & 2,5 & 0,25 & 0,16 & 0,00 & 976 \\
D10c & 700 & 45 & 40 & 2,0 & 7.000 & 0,41 & 0,8 & 0,15 & 0,09 & 0,00 & 2.070 \\
\bottomrule
\end{tabular}
}
\end{anhangblock}

\section{Ausgaben}
\label{sec:anhang-ausgaben}

\begin{anhangblock}
\captionof{table}{Ausgaben des Gesamtstaats 2025 nach Posten (\datei{STAATSAUSGABEN}).
Bürger- und Kindergeld kommen aus den Haushaltsprofilen hinzu, ebenso die
Änderungen durch Renten, Demografie, Rechtsstand und Erhebungskosten
(Abschnitt~\ref{sec:ausgaben}). Der Ausgleichsposten bringt den Saldo 2025 auf
den Finanzierungssaldo der VGR.}
\label{tab:ausgaben}
{\small% Erzeugt von docs/modell/messung.js — nicht von Hand bearbeiten
\begin{tabular}{@{}lr@{}}
\toprule
Posten & Mrd.\,\euro{} \\
\midrule
Soziale Sicherung & 850,0 \\
Gesundheit & 341,0 \\
Bildung & 180,0 \\
Verteidigung & 90,0 \\
Infrastruktur & 120,0 \\
Allgemeine Verwaltung & 140,0 \\
Zinsen & 54,0 \\
Sonstiges & 140,0 \\
Grundsicherung ohne Regelsatz (Unterkunft, Mehrbedarfe) & 10,8 \\
Klima- und Transformationsfonds & 12,6 \\
Ausgleichsposten: nicht modellierte Einnahmen & $-$221,4 \\
\midrule
Summe & 1.717,0 \\
\bottomrule
\end{tabular}
}
\end{anhangblock}

\section{Ressortgraphen}
\label{sec:anhang-graphen}

\newcommand{\ressortgraph}[3]{%
\begin{anhangblock}
\begin{tikzpicture}
\input{generiert/graph_#1.tex}
\end{tikzpicture}
\captionof{figure}{#2}
\label{#3}
\end{anhangblock}}

Die Abbildungen~\ref{fig:finanzen} bis~\ref{fig:umwelt} zeigen die gemessenen
Kanten von den Stellgrößen am Tisch zu den Kanälen, erste Periode, je Schritt.
Links stehen die Stellgröße, ihr Schritt und ihre Saldowirkung, rechts die
Kanäle mit ihrer Einheit, an jeder Kante die Änderung des Kanals. Blau heißt,
der Kanal steigt, orange, er sinkt. Die Strichstärke wächst mit dem Betrag,
bezogen auf den größten Wert derselben Einheit über alle vier Ressorts. Kanten
unter 0,3~Mrd.~\euro{} beziehungsweise 0,02\,\% entfallen;
Tabelle~\ref{tab:kanaele} enthält alle Werte.

\ressortgraph{Finanzen}{Ressort Finanzen: gemessene Wirkung auf die Kanäle,
erste Periode, je Schritt.}{fig:finanzen}
\ressortgraph{Wirtschaft}{Ressort Wirtschaft. Die Investitionsreaktion wirkt
zusätzlich über das private Kapital auf das BIP der Folgeperioden
(Abbildung~\ref{fig:struktur}).}{fig:wirtschaft}
\ressortgraph{Soziales}{Ressort Soziales. Beitragssätze wirken unmittelbar nur
auf die Einnahmen, das Rentenniveau nur auf die Ausgaben; alle übrigen Kanten
laufen über Konsum und Nachfrage.}{fig:soziales}
\ressortgraph{Umwelt}{Ressort Umwelt. Das Klimageld verschiebt 70\,\% des
\COz-Aufkommens zu den Transfers.}{fig:umwelt}

\section{Messwerte}
\label{sec:anhang-messwerte}

\begin{anhangblock}
\captionof{table}{Wirkung auf die Kanäle, erste Periode (2025--2028, Jahreswert), je
Schritt. Das Klimageld erscheint bei den Transfers und als Minderung der
\COz-Einnahmen, als die der Staat es bucht; in den Saldo geht es einmal ein.}
\label{tab:kanaele}
{\scriptsize\setlength{\tabcolsep}{2.6pt}% Erzeugt von docs/modell/messung.js — nicht von Hand bearbeiten
\begin{tabular}{@{}l*{11}{r}@{}}
\toprule
 & \multicolumn{6}{c}{Einnahmen (Mrd.\,\euro{})} & \multicolumn{2}{c}{Ausgaben} & \multicolumn{2}{c}{Verhalten (\%)} & Nachfr. \\
\cmidrule(lr){2-7}\cmidrule(lr){8-9}\cmidrule(lr){10-11}
Stellgröße (Schritt) & ESt & USt & Untern. & Verm. & Beitr. & CO\textsubscript{2} & Transf. & Renten & Arbeit & Invest. & Mrd.\,\euro{} \\
\midrule
Grundfreibetrag (+1.000\,\euro{}) & $-$12,6 & $+$3,4 & $+$0,2 & $+$0,1 & $+$1,4 & 0,0 & 0,0 & 0,0 & $+$0,35 & 0,00 & $+$8,5 \\
Eingangssteuersatz (+1\,Pp.) & $+$4,5 & $-$1,1 & $-$0,1 & 0,0 & $-$0,5 & 0,0 & 0,0 & 0,0 & $-$0,08 & 0,00 & $-$2,8 \\
Spitzensteuersatz (+1\,Pp.) & $+$0,2 & 0,0 & 0,0 & 0,0 & 0,0 & 0,0 & 0,0 & 0,0 & 0,00 & 0,00 & $-$0,1 \\
Umsatzsteuer (+1\,Pp.) & $-$0,5 & $+$9,5 & $-$0,1 & 0,0 & $-$0,9 & 0,0 & 0,0 & 0,0 & 0,00 & 0,00 & $-$5,5 \\
Erbschaftsteuer (+1\,Pp.) & 0,0 & 0,0 & 0,0 & $+$0,3 & 0,0 & 0,0 & 0,0 & 0,0 & 0,00 & 0,00 & $-$0,1 \\
Körperschaftsteuer (+1\,Pp.) & $-$0,1 & $-$0,1 & $+$2,5 & 0,0 & $-$0,2 & 0,0 & 0,0 & 0,0 & 0,00 & $-$0,40 & $-$1,0 \\
Gewerbesteuer (+1\,Pp.) & $-$0,2 & $-$0,1 & $+$4,8 & 0,0 & $-$0,3 & 0,0 & 0,0 & 0,0 & 0,00 & $-$0,40 & $-$2,0 \\
RV-Beitrag (+1\,Pp.) & $-$0,9 & $-$2,3 & $-$0,3 & $-$0,1 & $+$15,7 & 0,0 & 0,0 & 0,0 & 0,00 & 0,00 & $-$10,3 \\
KV-Beitrag (+1\,Pp.) & $-$0,8 & $-$2,1 & $-$0,3 & $-$0,1 & $+$15,9 & 0,0 & 0,0 & 0,0 & 0,00 & 0,00 & $-$9,6 \\
Rentenniveau (+1\,Pp.) & $+$0,4 & $+$0,3 & $+$0,1 & 0,0 & $+$0,8 & 0,0 & 0,0 & $+$7,5 & 0,00 & 0,00 & $+$4,8 \\
Bürgergeld (+10\,\euro{}/Monat) & 0,0 & 0,0 & 0,0 & 0,0 & $+$0,1 & 0,0 & $+$0,5 & 0,0 & 0,00 & 0,00 & $+$0,4 \\
Kindergeld (+10\,\euro{}/Monat) & $+$0,1 & $+$0,1 & 0,0 & 0,0 & $+$0,2 & 0,0 & $+$2,0 & 0,0 & 0,00 & 0,00 & $+$1,3 \\
CO\textsubscript{2}-Preis (+10\,\euro{}/t) & $-$0,1 & $-$0,1 & 0,0 & 0,0 & $-$0,2 & $+$2,6 & 0,0 & 0,0 & 0,00 & 0,00 & $-$1,4 \\
Klimageld (ein) & $+$0,6 & $+$0,5 & $+$0,2 & $+$0,1 & $+$1,3 & $-$12,6 & $+$12,6 & 0,0 & 0,00 & 0,00 & $+$7,8 \\
\addlinespace
\multicolumn{12}{@{}l}{\emph{Weitere Stellgrößen}} \\
Grenze Spitzensatz ($-$10.000\,\euro{}) & 0,0 & 0,0 & 0,0 & 0,0 & 0,0 & 0,0 & 0,0 & 0,0 & 0,00 & 0,00 & 0,0 \\
Umsatzsteuer, ermäßigt (+1\,Pp.) & $-$0,2 & $+$5,3 & $-$0,1 & 0,0 & $-$0,5 & 0,0 & 0,0 & 0,0 & 0,00 & 0,00 & $-$2,9 \\
Abgeltungsteuer (+1\,Pp.) & $+$2,8 & $-$0,2 & 0,0 & 0,0 & $-$0,1 & 0,0 & 0,0 & 0,0 & 0,00 & 0,00 & $-$0,8 \\
Beitragsbemessungsgrenze (+10.000\,\euro{}) & $-$0,3 & $-$1,5 & $-$0,1 & 0,0 & $+$8,3 & 0,0 & 0,0 & 0,0 & 0,00 & 0,00 & $-$3,0 \\
Vermögensteuer (+0,1\,Pp.) & $-$0,1 & $-$0,1 & 0,0 & $+$3,4 & $-$0,2 & 0,0 & 0,0 & 0,0 & 0,00 & 0,00 & $-$1,3 \\
Bodenwertsteuer (+0,1\,Pp.) & $-$0,2 & $-$0,1 & $-$0,1 & $+$5,0 & $-$0,3 & 0,0 & 0,0 & 0,0 & 0,00 & 0,00 & $-$2,0 \\
Mindeststeuer Milliardäre (+1\,Pp.) & $-$0,1 & 0,0 & 0,0 & $+$3,9 & $-$0,1 & 0,0 & 0,0 & 0,0 & 0,00 & 0,00 & $-$0,7 \\
\bottomrule
\end{tabular}
}
\end{anhangblock}

\begin{anhangblock}
\captionof{table}{Wirkung auf die Kennzahlen, erste Periode (2025--2028, Jahreswert), je
Schritt. Gini in Punkten, Nettoeinkommen je Haushalt und Jahr.}
\label{tab:kennzahlen}
{\footnotesize\setlength{\tabcolsep}{3pt}% Erzeugt von docs/modell/messung.js — nicht von Hand bearbeiten
\begin{tabular}{@{}l*{8}{r}@{}}
\toprule
 & Saldo & Gini & Armut & Emiss. & BIP & \multicolumn{3}{c@{}}{Nettoeinkommen (\euro{})} \\
\cmidrule(l){7-9}
Stellgröße (Schritt) & Mrd.\,\euro{} & Pkt. & Pp. & Mt & Mrd.\,\euro{} & D1 & D5 & D10c \\
\midrule
Grundfreibetrag (+1.000\,\euro{}) & $-$6,9 & $-$0,123 & $+$0,29 & 0,0 & $+$8,5 & 0 & $+$357 & $+$637 \\
Eingangssteuersatz (+1\,Pp.) & $+$2,6 & $-$0,036 & $-$0,09 & 0,0 & $-$2,8 & 0 & $-$100 & $-$322 \\
Spitzensteuersatz (+1\,Pp.) & $+$0,1 & $-$0,024 & 0,00 & 0,0 & $-$0,1 & 0 & 0 & $-$1.609 \\
Umsatzsteuer (+1\,Pp.) & $+$7,7 & $+$0,057 & $-$0,03 & 0,0 & $-$5,5 & $-$87 & $-$204 & $-$1.555 \\
Erbschaftsteuer (+1\,Pp.) & $+$0,3 & $-$0,007 & 0,00 & 0,0 & $-$0,1 & 0 & $-$2 & $-$177 \\
Körperschaftsteuer (+1\,Pp.) & $+$2,0 & $-$0,035 & $-$0,01 & 0,0 & $-$1,0 & $-$8 & $-$27 & $-$1.616 \\
Gewerbesteuer (+1\,Pp.) & $+$3,9 & $-$0,068 & $-$0,02 & 0,0 & $-$2,0 & $-$16 & $-$52 & $-$3.128 \\
RV-Beitrag (+1\,Pp.) & $+$11,7 & $+$0,064 & $-$0,10 & 0,0 & $-$10,3 & $-$129 & $-$366 & $-$987 \\
KV-Beitrag (+1\,Pp.) & $+$12,1 & $+$0,119 & $-$0,10 & 0,0 & $-$9,6 & $-$129 & $-$366 & $-$679 \\
Rentenniveau (+1\,Pp.) & $-$5,8 & $-$0,216 & $-$0,01 & 0,0 & $+$4,8 & $+$122 & $+$238 & $+$117 \\
Bürgergeld (+10\,\euro{}/Monat) & $-$0,4 & $-$0,042 & $-$0,08 & 0,0 & $+$0,4 & $+$72 & 0 & 0 \\
Kindergeld (+10\,\euro{}/Monat) & $-$1,7 & $-$0,066 & $-$0,02 & 0,0 & $+$1,3 & $+$45 & $+$59 & $+$11 \\
CO\textsubscript{2}-Preis (+10\,\euro{}/t) & $+$2,0 & $+$0,019 & $-$0,01 & $-$9,8 & $-$1,4 & $-$24 & $-$55 & $-$303 \\
Klimageld (ein) & $-$9,7 & $-$0,317 & $-$0,15 & 0,0 & $+$7,8 & $+$307 & $+$307 & $+$307 \\
\addlinespace
\multicolumn{9}{@{}l}{\emph{Weitere Stellgrößen}} \\
Grenze Spitzensatz ($-$10.000\,\euro{}) & 0,0 & $-$0,001 & 0,00 & 0,0 & 0,0 & 0 & 0 & $-$72 \\
Umsatzsteuer, ermäßigt (+1\,Pp.) & $+$4,4 & $+$0,030 & $-$0,02 & 0,0 & $-$2,9 & $-$46 & $-$108 & $-$824 \\
Abgeltungsteuer (+1\,Pp.) & $+$2,3 & $-$0,072 & $-$0,01 & 0,0 & $-$0,8 & $-$1 & $-$11 & $-$2.984 \\
Beitragsbemessungsgrenze (+10.000\,\euro{}) & $+$6,2 & $-$0,229 & 0,00 & 0,0 & $-$3,0 & 0 & 0 & $-$1.431 \\
Vermögensteuer (+0,1\,Pp.) & $+$2,8 & $-$0,075 & $-$0,02 & 0,0 & $-$1,3 & 0 & $-$19 & $-$1.868 \\
Bodenwertsteuer (+0,1\,Pp.) & $+$4,2 & $-$0,109 & $-$0,02 & 0,0 & $-$2,0 & 0 & $-$27 & $-$2.726 \\
Mindeststeuer Milliardäre (+1\,Pp.) & $+$3,4 & $-$0,142 & 0,00 & 0,0 & $-$0,7 & 0 & 0 & $-$9.476 \\
\bottomrule
\end{tabular}
}
\end{anhangblock}

\begin{anhangblock}
\captionof{table}{Wirkung am Ende des Pfads gegenüber dem Status-quo-Pfad, die Stellgröße
über alle Perioden verschoben: Schuldenquote und BIP-Niveau zu Beginn der
letzten Periode (2041), \COz{} kumuliert 2025--2040, Saldo und Gini der letzten
Periode (2041--2044). Status quo: Schuldenquote \SQschuldEnde\,\%, BIP-Niveau
\SQbipEnde~Mrd.~\euro, \COz{} kumuliert \SQcoEnde~Mt.}
\label{tab:pfad}
{\footnotesize% Erzeugt von docs/modell/messung.js — nicht von Hand bearbeiten
\begin{tabular}{lrrrrr}
\toprule
Stellgröße (Schritt) & \rotatebox{75}{Schuldenquote 2041 [Pp.]} & \rotatebox{75}{BIP-Niveau 2041 [Mrd. \euro{}]} & \rotatebox{75}{CO\textsubscript{2} 2025–2040 [Mt]} & \rotatebox{75}{Saldo 2041–44 [Mrd. \euro{}/a]} & \rotatebox{75}{Gini 2041–44 [Gini-Pkt.]} \\
\midrule
Grundfreibetrag (1.000 \euro{}) & $+$0,85 & $+$15,3 & 0 & $-$3,6 & $-$0,124 \\
Eingangssteuersatz (1 Pp.) & $-$0,51 & $-$3,7 & 0 & $+$2,7 & $-$0,037 \\
Spitzensteuersatz (1 Pp.) & $-$0,04 & $-$0,1 & 0 & $+$0,2 & $-$0,027 \\
Umsatzsteuer (1 Pp.) & $-$2,44 & 0,0 & 0 & $+$14,7 & $+$0,057 \\
Erbschaftsteuer (1 Pp.) & $-$0,08 & 0,0 & 0 & $+$0,5 & $-$0,007 \\
Körperschaftsteuer (1 Pp.) & $-$0,33 & $-$6,5 & 0 & $+$1,4 & $-$0,038 \\
Gewerbesteuer (1 Pp.) & $-$0,93 & $-$6,5 & 0 & $+$5,0 & $-$0,072 \\
RV-Beitrag (1 Pp.) & $-$3,76 & 0,0 & 0 & $+$22,8 & $+$0,066 \\
KV-Beitrag (1 Pp.) & $-$3,90 & 0,0 & 0 & $+$23,6 & $+$0,122 \\
Rentenniveau (1 Pp.) & $+$2,01 & 0,0 & 0 & $-$13,3 & $-$0,275 \\
Bürgergeld (10 \euro{}/Monat) & $+$0,11 & 0,0 & 0 & $-$0,7 & $-$0,043 \\
Kindergeld (10 \euro{}/Monat) & $+$0,52 & 0,0 & 0 & $-$3,1 & $-$0,067 \\
CO\textsubscript{2}-Preis (10 \euro{}/t) & $-$0,46 & 0,0 & $-$103 & $+$1,7 & $+$0,007 \\
Klimageld (ein) & $+$2,19 & 0,0 & 0 & $-$8,2 & $-$0,119 \\
\bottomrule
\end{tabular}
}
\end{anhangblock}

\Needspace*{12\baselineskip}
\section{Elastizitäten}
\label{sec:anhang-elast}

Die Objekte \datei{ELAST} und \datei{ELAST\_QUELLEN} in \datei{js/data.js}. Die
Schlüssel \datei{capital\_supply} und \datei{evasion} stehen im Datensatz, das
Modell verwendet sie nicht.

{\small% Erzeugt von docs/modell/messung.js — nicht von Hand bearbeiten
\begin{longtable}{@{}p{2.6cm}rp{2.2cm}p{8.4cm}@{}}
\toprule
Schlüssel & Wert & Spanne & Fundstelle \\
\midrule
\endhead
\texttt{labor\_supply} & 0,200 & 0,1–0,3 & Saez/Chetty/Gruber Konsens · ifo Schnelldienst 01/2025 \\
\texttt{capital\_supply} & 0,500 & 0,4–0,8 & Kleven/Schultz (2014) AEJ: Economic Policy 6(4) \\
\texttt{consumption} & $-$0,350 & −0,2 bis −0,5 & Lewbel/Pendakur (2009) AER 99(3) · Metaanalyse Havranek et al. 2018 \\
\texttt{co2} & $-$0,300 & −0,2 bis −0,4 & EWI/DIW BEHG-Evaluation 2023 · Edenhofer/PIK 2024 \\
\texttt{evasion} & 0,250 & 0,1–0,3 & Schneider (2023) Shadow Economy DE · IfW Kiel 2024 \\
\texttt{investment} & $-$0,400 & −0,3 bis −0,5 & Gechert/Heimberger (2022) European Economic Review 147 · Neumeier SVR Arbeitspapier 03/2025 \\
\texttt{d10c\_labor} & 0,400 & 0,25–0,5 & Saez/Slemrod/Giertz (2012) JEL 50(1) · Piketty/Saez/Stantcheva (2014) AEJ: Economic Policy 6(1) \\
\texttt{d10c\_avoidance} & 0,500 & 0,3–0,7 & Kleven/Schultz (2014) AEJ: Economic Policy 6(4) · Chetty/Friedman/Saez (2013) \\
\texttt{d10c\_wegzug} & 0,100 & 0,05–0,15 & Beleg ausstehend \\
\texttt{verm\_ausweichen} & 0,215 & 0,1–0,43 je Pp. & Brülhart/Gruber/Krapf/Schmidheiny (2022) AEJ: Economic Policy 14(4), Behavioral Responses to Wealth Taxes \\
\bottomrule
\end{longtable}
}

\Needspace*{12\baselineskip}
\section{Formelquellen}
\label{sec:anhang-quellen}

Die Objekte \datei{FORMEL\_QUELLEN\_*} der Rechenmodule, unverändert
abgedruckt.

{\small% Erzeugt von docs/modell/messung.js — nicht von Hand bearbeiten
\begin{longtable}{@{}p{4.1cm}p{11.4cm}@{}}
\toprule
Quelle & Formel und Fundstelle \\
\midrule
\endhead
\multicolumn{2}{@{}l}{\textbf{\texttt{berechne.js}}} \\
\texttt{arbeitsangebot} & {\formelschrift labor\_factor = 1 + ε × Δ(1 − GS) / (1 − GS\_SQ)}\newline Saez/Chetty/Gruber Konsens · ifo Schnelldienst 01/2025 · Piketty/Saez/Stantcheva (2014) AEJ: Economic Policy 6(1) \\
\texttt{bge\_arbeitsangebot} & {\formelschrift lf = lf\_tax × (1 − BGE\_LABOR\_EFF[i] × min(1,67; BGE/1200))}\newline RWI (2024) bis −30 \% bei 1.500 \euro{} · DIW Pilot 2024 (n=107) · ZEW Heim et al. (extensiver Margin) · ifo Mikrosimulation 2021 \\
\texttt{mwst\_konsumanteil} & {\formelschrift MwSt = Konsum × (0,70 × t\_reg/(1+t\_reg) + 0,30 × t\_erm/(1+t\_erm))}\newline Destatis VGR 2024 (ca. 70 \% Regelsatz-Konsum) · Lewbel/Pendakur (2009) AER 99(3) \\
\texttt{co2\_emissionen} & {\formelschrift Emissionen = Pfad(Jahr) × [(1 − s) + s × max(0,4; min(1,1; 1 + ε\_CO₂ × (p − 55)/100))],  s = 327/649}\newline BEHG § 10 · EWI/DIW BEHG-Evaluation 2023 · Edenhofer/PIK 2024 · UBA Projektionsbericht 2025 (Pfad) · UBA Treibhausgas-Emissionen 2025 \\
\texttt{erbschaft} & {\formelschrift Erb\_auf = Masse\_top × Satz\_eff + Masse\_unten × min(Satz,15\%)/2}\newline Bach/Sinclair/Bührle/Wichers DIW 2025 · §§ 10, 19 ErbStG · BRH 2016 \\
\texttt{zucman} & {\formelschrift Zucman\_auf = Milliardärsvermögen × max(0; Satz − 0,3 \% ESt − VermSt-Satz) × (1 − 0,15 × min(1; Satz/2))}\newline Zucman (2024) A blueprint for a coordinated minimum effective taxation standard for ultra-high-net-worth individuals, G20-Bericht · EU Tax Observatory 2024 · Jakobsen/Kleven/Kolsrud NBER 2024 \\
\texttt{sv\_beitraege} & {\formelschrift SV = Lohnsumme\_sv × Satz\%  (nur bis BBG)}\newline § 158 SGB VI · § 241 SGB V · § 341 SGB III · § 55 SGB XI · DRV Beitragssätze 2025 \\
\texttt{rv\_ausgaben} & {\formelschrift ΔRente\_i = Brutto\_i × rente\_anteil\_i × renten\_faktor × (Rentenniveau/48 \% − 1);  ΔRV-Ausgaben = Σ Haushalte ΔRente\_i;  KV, AL, PV fest}\newline § 68 SGB VI (Rentenanpassung) · § 154 Abs. 3 SGB VI (Sicherungsniveau) · Rentenpaket 2025 (BMAS, Haltelinie 48 \% bis 2031) · § 158 SGB VI · BMAS Rentenversicherungsbericht 2025 \\
\texttt{verwaltungskosten} & {\formelschrift Admin = ∑ Aufkommen\_i × Quote\_i  (ADMIN\_QUOTE je Steuer)}\newline BRH Jahresberichte · OECD Tax Administration 2024 · Normenkontrollrat Jahresbericht 2024 \\
\texttt{deadweight\_loss} & {\formelschrift DWL = ∑ᵢ 0,5 × ε × GS\_i² / (1 − GS\_i) × Lohnsumme\_i  (Harberger-Dreieck je Dezil)}\newline Harberger (1964) · Saez (2001) Review of Economic Studies 68(1) · Chetty (2009) AEJ: Economic Policy 1(2) \\
\texttt{dynamisches\_scoring} & {\formelschrift Δ\_dyn = Δ\_KSt × (investment\_factor − 1) + Δ\_ESt × (avg\_labor − 1)}\newline CBO Dynamic Scoring Guidelines · ifo Schnelldienst 01/2025 · Saez/Chetty Konsens \\
\texttt{inzidenz} & {\formelschrift Last\_i = ΔKSt+ΔGewSt × (½ Anteil Arbeitseinkommen\_i + ½ Anteil Kapitaleinkommen\_i) + ΔErbSt+ΔVermSt+ΔBodenSt × Anteil Vermögen\_i + ΔZucman × D10c}\newline Fuest/Peichl/Siegloch (2018) AER 108(2): Beschäftigte tragen rund 51 \% der Gewerbesteuer · Harberger (1962) · CBO (2012) Distribution of Corporate Tax · Mirrlees Review (2011) Tax by Design, Kap. 16 (Bodenwertsteuer) \\
\midrule
\multicolumn{2}{@{}l}{\textbf{\texttt{einkommensteuer.js}}} \\
\texttt{estTarif} & {\formelschrift ∫₀ˣ r(z) dz  (stückweise lineare Grenzsteuerrate, 5 Zonen)}\newline § 32a Abs. 1 EStG 2026 (Steuerfortentwicklungsgesetz, BGBl. 2024 I Nr. 449) · Formeltarif (kontinuierlich, keine Sprünge) \\
\texttt{grenzsteuersatz} & {\formelschrift r(z) = r₀ + (rₘ − r₀) × z/g₂  [Zone 2], linear interpoliert je Zone}\newline § 32a Abs. 1 Nr. 2–5 EStG · BMF Steuerschätzung 2025 \\
\texttt{effSteuersatz} & {\formelschrift T(z) / z}\newline § 2 Abs. 5 EStG \\
\texttt{spitzenzone} & {\formelschrift Zuschlag\_D10c = (r5 − r4) × E[max(0, x − Grenze)],  x \textasciitilde{} Pareto(α = 1,5) mit dem zvE-Mittel des Dezils}\newline Atkinson/Piketty/Saez (2011) JEL 49(1), Top Incomes in the Long Run of History · Bartels/Jenderny (2015) World Top Incomes Database Working Paper 2015/1 \\
\texttt{estHaushalt} & {\formelschrift T\_HH = s × T(brutto × q / s)  mit q = zvE-Quote, s = Splitting-Faktor}\newline §§ 26, 32a Abs. 5 EStG (Splitting) · Destatis ESt-Statistik (zvE vs. Bruttoeinkünfte) · Mikrozensus 2024 \\
\midrule
\multicolumn{2}{@{}l}{\textbf{\texttt{verteilung.js}}} \\
\texttt{aequivalenzEinkommen} & {\formelschrift y\_äq,i = Netto\_i / Bedarfsgewicht\_i  (neue OECD-Skala: 1,0 erste Person · 0,5 weitere ab 14 J. · 0,3 Kinder unter 14 J.)}\newline Eurostat EU-SILC Methodik (äquivalisiertes verfügbares Einkommen) · Destatis Glossar „Äquivalenzeinkommen" · OECD (2013) Framework for Statistics on the Distribution of Household Income \\
\texttt{berechneGini} & {\formelschrift G = 1 − 2·∫Lorenz(x)dx  (Trapezregel, gewichtet nach Haushaltszahl)}\newline Sen (1973) On Economic Inequality · Cowell (2011) Measuring Inequality · Destatis Methodik Gini-Koeffizient \\
\texttt{berechnePalma} & {\formelschrift Palma = Einkommensanteil(Top 10\%) / Einkommensanteil(Bottom 40\%)}\newline Palma (2011) Homogeneous Middles vs. Heterogeneous Tails · UNDP HDR 2013 \\
\texttt{berechneMedianGewichtet} & {\formelschrift Gewichteter Median: kumulierte Haushaltsanteile bis 50 \%}\newline Destatis Mikrozensus 2024 · SOEP v40 \\
\texttt{armutsrisiko} & {\formelschrift pov\_i = pov\_sq\_i × (y\_i/y\_sq\_i ÷ PL/PL\_sq)\textasciicircum{}(−1,5)}\newline Bourguignon (2003) · EU-SILC DE 2024: 15,5 \% (Destatis) · SOEP v40 · IAB Kurzbericht 2024 \\
\texttt{berechneNettoSQ} & {\formelschrift Netto\_SQ = Brutto − ESt(SQ) − SV(SQ) − MwSt(SQ) − CO₂(SQ) + Klimageld(SQ) + Transfers(SQ)}\newline PRESETS.status\_quo (data.js) · § 32a EStG 2025 · § 158 SGB VI · § 241 SGB V \\
\texttt{berechneDezilDelta} & {\formelschrift Δ\_i = Netto\_neu\_i − Netto\_SQ\_i}\newline SOEP v40 · Destatis Mikrozensus 2024 · BMF Steuerschätzung 2025 \\
\bottomrule
\end{longtable}
}
